\begin{center}
    {\LARGE \textbf{2003无机化学}} \\[2ex]
    {\large 时量180分钟 \quad 满分100分}
\end{center}

\vspace{0.5cm}
\noindent\textbf{注意事项:}
\begin{enumerate}[label=\arabic*., parsep=0pt, itemsep=0pt]
    \item 答题前,考生先将自己的姓名、学号、考场号、座位号填写在答题卡上,将条形码准确粘贴在考生信息条形码粘贴区。
    \item 考试结束后,将本试卷与答题纸一并收回。
    \item 回答各个问题时,将答案写在答题纸上,写在本试卷上无效。
\end{enumerate}

\vspace{0.5cm}
\noindent\textbf{一、按要求回答下列问题。}

\begin{enumerate}[label=\arabic*., start=1, itemsep=1.5ex]
    \item 完成并配平下列化学反应方程式(10 pts.):
    \begin{enumerate}[label=(\arabic*), itemsep=1ex]
        \item 酸性\ce{KMnO4}溶液与\ce{NaNO2}溶液反应。
        \item 重铬酸铵的热分解
        \item 金属铂溶于王水
        \item 向\ce{FeCl3}溶液中通入\ce{H2S}气体
        \item \ce{H2O}滴加在红磷和碘的混合物上
    \end{enumerate}

    \item 按顺序写出原子序数从21到30的10种元素的元素名称和元素符号。在元素名称下划横线表示出具有5个d电子的元素，在元素符号下划横线表示出电子构型“特殊”的元素。(16 pts.)

    \item 根据价层电子对互斥理论和杂化轨道理论，给出下列分子和离子的价层电子对数，电子对的空间构型，分子或离子的几何构型和中心原子的杂化方式。(16 pts.)

    {\centering
    \ce{CO2}，\ce{NF3}，\ce{SF4}，\ce{I3^-}，\ce{H3O+} \par
    }

    \item 写出\ce{H3BO3}、\ce{H3PO3}在水中显酸性的反应方程式；写出\ce{NH3}、\ce{NH2OH}在水中显碱性的反应方程式，并说明\ce{NH2OH}比\ce{NH3}的碱性弱的原因。(16 pts.)

    \item 命名\ce{[Cr(H2O)4Br2]Br\cdot 2H2O}，其内界有几种异构体?试用简图分别表示出其结构。画出中心离子的d电子在晶体场中的排列方式，并求出其晶体场稳定化能。(16 pts.)

    \item 已知：
        
        $E^\ominus(\ce{Co^{3+}/Co^{2+}}) = 1.81\ \text{V}$

        $E^\ominus(\ce{[Co(NH3)6]^{3+}/[Co(NH3)6]^{2+}}) = 0.04\ \text{V}$

        试指出$K'_{\text{稳}}(\ce{[Co(NH3)6]^{3+}})$和$K''_{\text{稳}}(\ce{[Co(NH3)6]^{2+}})$哪一个大些，并求出$K'/K''$的值。(16 pts.)

\end{enumerate}

\vspace{0.5cm}
\noindent\textbf{二、制备、分离与鉴别。}

\begin{enumerate}[label=\arabic*., start=1, itemsep=1.5ex]
    \item 以\ce{NH4HSO4}为主要原料，利用电解-水解法生产\ce{H2O2}，是20世纪初开发的化工项目。试用化学方程式表示这一生产过程。(10 pts.):

    \item 设计方案分离下列混合离子：\ce{Ag^{+}},\ce{Hg^{2+}},\ce{Co^{2+}},\ce{Al^{3+}},\ce{Fe^{3+}},\ce{Cr^{3+}}。(10 pts.)

    \item 三种白色固体分别是\ce{KClO},\ce{KClO3},\ce{KClO4}，试加以鉴别。(10 pts.)

\end{enumerate}

\vspace{0.5cm}
\noindent\textbf{三、综合题}

\begin{enumerate}[label=\arabic*., start=1, itemsep=1.5ex]
    \item 试用离子极化学说分析和解释下列化合物分解温度数据中所呈现的某些规律性。(10 pts.)
    
    \vspace{1ex}
    \begin{tabular}{l r @{\qquad} l r @{\qquad} l r}
        \ce{MgSO4} & 895$^\circ$C &
        \ce{CaSO4} & 1149$^\circ$C &
        \ce{SrSO4} & 1374$^\circ$C \\

                   &              &
        \ce{CdSO4} & 816$^\circ$C &
        \ce{MnSO4} & 755$^\circ$C \\

                   &              &
                   &              &
        \ce{Mn2(SO4)3} & 300$^\circ$C \\
    \end{tabular}
    \vspace{1ex}
    
    并对\ce{BaSO4}、\ce{ZnSO4}的分解温度做出估计。

    \item 有一白色固体，可能是\ce{KI}、\ce{CaI2}、\ce{KIO3}和\ce{BaCl2}中的一种或两种的混合物。试根据下述实验判断白色固体的组成。(10 pts.)
    \begin{enumerate}[label=(\arabic*), itemsep=1ex]
        \item 将白色固体溶于水得无色溶液。
        \item 向此溶液中加入少量的稀硫酸，溶液变黄并有白色沉淀生成，溶液遇淀粉立刻变蓝。
        \item 向蓝色溶液中加入\ce{NaOH}溶液后，蓝色消失而白色沉淀并未消失。
    \end{enumerate}

    \item 讨论第二周期元素的氢化物的结构和性质。(10 pts.)
\end{enumerate}

